\section{Заключение}
Часть задач, поставленных в начале исследования, были выполнены.
Некоторые задачи потребовали гораздо больше усилий, чем изначально
предполагалось. Возникли неожиданные сложности в реализации, но
в итоге основной результат был достигнут. Нейросеть Хопфилда
была обучена алгоритмом скоростного градиента и результат оказался
лучше, чем при обучении по классическому алгоритму Хебба.

\textbf{Выводы:}
В ходе работы была исследована применимость скоростного градиента к
обучению классической непрерывной сети Хопфилда.

Были выполнены почти все задачи исследования:
\begin{itemize}
  \item
        Была реализована классическая и непрерывная модели Хопфилда;
  \item
        Было проведено обучение с помощью правила Хэбба, Сторки и скоростного
        градиента;
  \item
        Продемонстрированы результаты с разными функциями шума и
        параметром регуляризации.
\end{itemize}


Следующими задачами исследования являются:
\begin{itemize}
  \item Рассмотрение нескольких других методов решения;
  \item Ввод метрик для точного сравнения;
  \item Приведение доказательств основных фактов сходимости, устойчивости и других;
  \item Применимость к задаче LMAPF.
\end{itemize}